\section*{Limitations}

\paragraph{Models.} We test seven generators from three families: four OpenAI models (gpt-4o-mini, gpt-4.1 and two
current reasoning models), Claude Haiku 5.5 and Sonnet 5.5, and Qwen2.5-7B at 4-bit precision. Gemini serves only as a judge, and we run
no large open model. Qwen2.5-7B, gpt-6.1-terra and Claude Sonnet ran with one seed, as did gpt-4.1 outside RAMDocs, and the reasoning and Claude models
only at their default temperature, so their sampling is not deterministic, and Claude ran through Anthropic's
OpenAI-compatible endpoint. The crossed reliability test (H24) covers \hTwentyFourCells\ detector--generator cells on one benchmark, with six
detectors, so any test of a detector-level property has little power (an exact permutation has \hRedoNPerm\
arrangements); the registered gap turned out to be dominated by the generator's own accuracy. The follow-rate
account of the crossed cells was defined after seeing them and uses follow rates measured on the same questions, so it
is a retrospective decomposition; no outcome-blind question set tests it, and the break-even precision is a
calibration for randomised flags on true pairs, not a validated rule for selective detectors. Concurrent runs gave
a detector slightly different outputs in different cells; its output is identical across its cells for at least
\frozenShareMin\% of questions, and restricting each cell to those moves no cell's gain by more than \frozenMaxShift\
points. The LLM detector is a single
prompt. The published baselines are reimplemented (Astute RAG has no official code; restoring TrustRAG's system message
changed little), NLI sees only a sentence-sized window, and we do not compare
against a trained detector such as ConflictRAG's \citep{wang2026conflictrag}.

\paragraph{Data and power.} The confirmatory evidence comes from two benchmarks: RAMDocs (300 questions, 139 of them
conflict questions) and ConflictBank (400 sampled questions, 300 with conflicts). In both, the misinformation is
synthetic, made by editing or generating text around a substituted answer; ConflictBank's passages are LLM-written
on both sides. RAGuard's misleading side is real text, but its labels are behavioural (whether GPT-4 judged the claim
correctly with the post), and an LLM audit finds only \rgAudContradict\% of sampled supporting--misleading pairs to
contradict (Appendix~\ref{app:semantic}), so RAGuard's oracle is a label transfer, not a factual oracle. With its posts in context
generators abstain on \rgAbstMin--\rgAbstMax\% of conflict claims, which is why accuracy on this true/false task
falls below chance; on the claims they answer it is \rgAnsMin--\rgAnsMax\%, against \rgCbMin--\rgCbMax\% closed-book. RAMDocs' document labels are imperfect. We exclude 35 hold-out
questions whose status depends on documents we could not label without human annotation, and we check that the
conclusions survive treating those documents as misinformation. Our comparisons detect effects of about
\mdeMed\ points with 80\% power, so smaller effects, including small harms on clean questions, may be missed.
WikiContradict has genuine conflicts but no misinformation side. CONFLICTS shows no significant effect for any
condition; scored only on items with a reference answer, it can neither confirm nor refute our findings. The majority variant tests one case where the false side is
more numerous, but its extra wrong-side passages are generated for the same substitution, which a real attacker
would vary, and they include ConflictBank's temporal and semantic conflicts, which we count as misinformation. On
RAGuard every generator is more accurate closed-book than with the posts (gpt-4o-mini \zcRgMiniConflict\% against
\zpRgMiniConflict\%), so its posts make a weak retrieval setting. The current
reasoning models distrust ConflictBank's generated passages when asked to judge them (Appendix~\ref{app:more}),
which suggests that generated benchmarks may not measure how such models treat real documents. Reversing and
paraphrasing the passages change the detectors' hit rate insignificantly (H29), but equivalence within $\pm$5 points
is shown for only one of four changes, so position and surface style are not ruled out as sources of their direction,
nor are deeper regularities of generated text.

\paragraph{Metric.} RAMDocs ships no scorer. Our primary metric reconstructs the published definition and scores our
MADAM-RAG port close to the published number on the pilot split, but the published subset is unknown, so this is not a
reproduction. The metric is string matching: it misses paraphrases, and partial answers interact with ``no wrong
answer''. An outside-family LLM judge agrees with it on \semAgree\% of \semN\ paired outputs; the matcher's misses
depend on the condition (more under answer keys and directional flags) but move no contrast by more than \semDidMax\
points (Appendix~\ref{app:semantic}). That sample cannot confirm the $\pm$3-point equivalence claims, and no human
labels check either scorer: the paired, blinded audit sheet we release is unlabelled. We report a more lenient variant throughout and mark any conclusion that depends on the choice.

\paragraph{Direction and the grid.} Oracle direction is not distinguishable from an explicit answer key (H23), so the
oracle rows are ceilings, not findings about detectors. In the synthetic grid the gain rises with direction accuracy, but the grid cannot separate adjacent levels. The gpt-4o-mini LLM
detector's direction precision (of its flags on a true contradiction, the share naming the false side; distinct from
its pair precision, the share of its flagged pairs that are true contradictions) is \dirAccLlm\% on RAMDocs and
\cbDirAccLlm\% on ConflictBank; its hit rate per question is lower, because it also misses pairs. Its
directional flag's gain over all ConflictBank questions with gpt-4o-mini is not significant after correction (H15),
and with stronger generators it hurts. Synthetic false flags are spread uniformly and their suspect side is random, whereas
real detectors concentrate their errors. The grid predicts real detectors only roughly, and three detectors cannot
show whether document- or pair-level summaries predict better.

\paragraph{Pre-registration.} The original hypotheses (H1--H9) were registered before the confirmatory run. H10--H31
were registered in eight addenda, after the earlier results were known though before their own runs, and the only
timestamp is the repository's commit history, not a third-party registry. Holm controls error within each family, not
across the programme: of \fdrN\ registered tests, \fdrHolm\ are significant under Holm and \fdrBH\ under a
Benjamini--Hochberg FDR over all of them (both shown in the hypothesis tables). Several analyses were chosen after seeing
data and are labelled as such: the TOST margin, the single- and multi-answer split, the document-level placement
axes, the post-hoc comparison of detector and
generator reliability that H24 later tested, H24's mechanism score, and the refinement of permutation $p$-values at
their floor.

\paragraph{Scope.} We study flags passed to the generator as text. Learned interfaces such as soft prompts or
credibility fine-tuning \citep{ye2026tcr,pan2024credibility}, and retrieval-side filtering, may follow different
trade-offs. We test one adversarial case, a detector that always names the wrong side, but not an attacker who controls
which pairs are flagged.

\section*{Ethics Statement}
This work uses public benchmarks and commercial model APIs, and involves no human subjects or personal data. Its main
risk is dual use. Our results show that a false ``likely misinformation'' label, or the removal of a passage, can
change a RAG system's answers: an inverted oracle that names the correct side as misinformation costs \aadvDirConf\
points on RAMDocs and \azCbMiniOracleInvFlagDirectionalConflict\ on ConflictBank as a flag, and \aadvRemConf\ on RAMDocs
when the passage it names is removed. An adversary who can trigger a deployed detector could exploit this
\citep{zou2025poisonedrag}. We therefore report harm on conflict-free questions alongside every gain, and we argue that
systems should measure what a detector does at the document level before acting on its flags.
